\clearpage\fancyhead[L]{\small PALLAS.}\hypertarget{native-327-lower}{}\label{native-327-lower}
\newcommand{\ObservedDeclinationMark}{\hbox{\vrule width0pt height6.625000bp depth0pt\special{pdf:bcontent}\special{pdf:code q 0 g 0.333333 6.625000 m 0.333333 3.791667 l 0.166667 3.750000 l 0.125000 3.541667 l 0.000000 3.541667 l 0.000000 3.041667 l 0.125000 3.041667 l 0.166667 2.875000 l 0.333333 2.833333 l 0.333333 0.916667 l 0.166667 0.875000 l 0.166667 0.166667 l 0.291667 0.166667 l 0.333333 0.000000 l 0.833333 0.000000 l 0.875000 0.166667 l 1.041667 0.166667 l 1.041667 0.875000 l 0.833333 0.916667 l 0.833333 2.833333 l 1.041667 2.875000 l 1.083333 3.041667 l 3.000000 3.041667 l 3.041667 2.875000 l 3.541667 2.875000 l 3.583333 3.041667 l 3.916667 3.041667 l 3.916667 3.208333 l 4.083333 3.250000 l 4.083333 3.375000 l 4.250000 3.416667 l 4.250000 3.541667 l 4.083333 3.583333 l 4.083333 3.750000 l 1.416667 3.750000 l 1.375000 3.541667 l 1.083333 3.541667 l 1.041667 3.750000 l 0.833333 3.791667 l 0.833333 6.083333 l 1.041667 6.125000 l 1.041667 6.458333 l 0.875000 6.458333 l 0.833333 6.625000 l h f* Q}\special{pdf:econtent}\kern4.250000bp}}

\begin{center}\textbf{PALLAS.}\par\vspace{4pt}\hrule\vspace{4pt}\fontsize{10}{12}\selectfont Göttingische gelehrte Anzeigen. Stück 8 u. 9. Seite 73\,..\,77. 1811 Jan. 14.\par\vspace{4pt}\hrule\vspace{4pt}\end{center}

\par\indent Im 198. Stück unserer Gelehrten Anzeigen vom vergangenen Jahr [S. 61 d. B.] haben wir über eine Untersuchung von Professor \textsc{Gauss} zu den sechs bisher beobachteten Oppositionen der Pallas berichtet, die sich nicht mehr durch eine rein elliptische Bewegung darstellen lassen. Professor \textsc{Gauss} war überzeugt, dass die immer größeren Abweichungen eine Folge der Störungen sind, die Pallas besonders durch Jupiter erfährt. Er begann deshalb schon im Oktober des vergangenen Jahres, die von Jupiter verursachten Störungen für die Jahre 1803..1811 zu berechnen. Zum einen wollte er sehen, um wie viel besser sich die Oppositionen unter Berücksichtigung dieser Störungen in Einklang bringen ließen. Zum anderen wollte er damit die spätere Entwicklung der allgemeinen Theorie dieser Störungen vorbereiten. Die heikle und schwierige Arbeit führte zu einem Ergebnis, das so befriedigend war, wie man es kaum hätte erwarten dürfen. Alle sechs Oppositionen ließen sich sowohl in der Länge als auch in der Breite bis auf wenige Sekunden in Einklang bringen. Da Pallas gegen Ende des Jahres 1810 wieder sichtbar werden musste, schien es interessant, vorab zu bestimmen, \SourceLeaf{328}wie stark sich die aus rein elliptischen Elementen berechnete Ephemeride für diese nächste Erscheinung durch die verbesserte Theorie des Planeten ändern würde, die die Jupiterstörungen berücksichtigt. Das lag umso näher, als Jupiters Einfluss bei der heliozentrischen Konjunktion von 1809 besonders stark gewesen war und die Elemente erheblich verändert hatte. Vor allem hatten Exzentrizität und mittlere Bewegung stark abgenommen, die mittlere Entfernung hatte zugenommen und die Apsidenlinie hatte sich zurückverlagert. Die Rechnung ergab eine Änderung der Ephemeride, die weit größer war als in irgendeinem der vorhergehenden Jahre. Gegen Ende Oktober war die Rektaszension um \emph{drei Viertel Grad} kleiner und die südliche Deklination um \emph{sechs Minuten} größer als in der Ephemeride, die in der Monatl. Corresp. vom Oktober 1810 abgedruckt ist. Diese Unterschiede nahmen immer weiter zu, wie die folgende Übersicht zeigt.
\Needspace{15\baselineskip}
\begin{center}Korrektur der Pallas-Ephemeride wegen der Störungen:\end{center}
\begin{center}\fontsize{10}{12}\selectfont\setlength{\tabcolsep}{4pt}\renewcommand{\arraystretch}{1.15}\begin{tabular}{llrr|r|r}
 &  &  &  & Rektaszension & Deklination\\\hline
1810. & Okt. & 26 & $0^{\text{u}}$ & $-45'$ & $+6'$\\
 & Nov. & 7 & 12 & $-52$ & $+7$\\
 & --- & 20 & 0 & $-59$ & $+9$\\
 & Dez. & 2 & 12 & $-67$ & $+11$\\
 & --- & 15 & 0 & $-75$ & $+12$\\
 & --- & 27 & 12 & $-84$ & \ObservedDeclinationMark\,$13$\\
1811 & Jan. & 9 & 0 & $-94$ & $+14$\\
 & --- & 21 & 12 & $-103$ & $+14$\\
 & Feb. & 3 & 0 & $-112$ & $+14$\\
\end{tabular}\end{center}

\par\indent Ein so großer Unterschied war für die Suche nach Pallas von größter Bedeutung. Professor \textsc{Gauss} teilte daher mehreren auswärtigen Astronomen sofort diese Ergebnisse mit, da das ungewöhnlich ungünstige Wetter die Suche hier verhinderte. Um in dieser wichtigen Angelegenheit ganz sicherzugehen, hielt es Professor \textsc{Gauss} jedoch für angebracht, die gesamte Störungsrechnung zu wiederholen. Bei dieser umfangreichen Arbeit konnten sich hier und da leicht Fehler eingeschlichen haben. Außerdem waren von der zweiten Rechnung schon deshalb zuverlässigere Ergebnisse zu erwarten, weil ihr bereits die verbesserte Theorie zugrunde gelegt werden konnte. Es zeigten sich aber nur unbedeutende Unterschiede von wenigen Sekunden. Die sechs Oppositionen stimmten weiterhin ebenso gut überein, und die erforderliche Änderung der Ephemeride blieb bis auf Kleinigkeiten ebenso groß. Damit man die jetzige Übereinstimmung der sechs Oppositionen besser einschätzen und mit den großen Abweichungen bei rein elliptischen Elementen (siehe unsere Anz. 1810, S. 1971) [S. 63 d. B.] vergleichen kann, geben wir sie hier an.
\Needspace{10\baselineskip}
\begin{center}\fontsize{10}{12}\selectfont\setlength{\tabcolsep}{4pt}\renewcommand{\arraystretch}{1.15}\begin{tabular}{r|r|r}
 & \multicolumn{2}{c}{Unterschied der}\\
 & \shortstack{mittlere\\Länge} & \shortstack{heliozentr.\\Breite}\\\hline
1803 & $+7''3$ & $+1''6$\\
1804 & $-6.8$ & $-4.0$\\
1805 & $+1.5$ & $-4.1$\\
1807 & $-3.3$ & $-3.8$\\
1808 & $-3.9$ & $-16.9$\\
1809 & $+5.0$ & $+3.5$\\
\end{tabular}\end{center}

\SourceLeaf{329}\par\indent Die Breite für 1808 war bekanntlich schlecht beobachtet worden. Daraus darf man schließen, dass die Störungen, die die anderen Planeten, besonders Saturn und Mars, in der Bewegung der Pallas verursachen, gegenüber denen durch Jupiter ganz unbedeutend sind und vorerst noch ganz vernachlässigt werden dürfen.

Auch die auswärtigen Astronomen, denen diese Ergebnisse mitgeteilt worden waren, hatten ungünstiges Wetter. Nur Herr von \textsc{Lindenau} beobachtete am 19. Dez. an der Seeberger Sternwarte einen kleinen Stern am Passageinstrument.
\[\begin{gathered}16^{\text{u}}\;9^{\text{m}}\;7^{\text{s}}1\ \text{M. Z.}\quad\text{Rektasz. }150^\circ\;20'\;12''45\\\text{geschätzte Deklination }20^\circ\;58'\text{ südl.}\end{gathered}\]
{\rightskip=0pt plus 1em \noindent Das stimmte auf die Minute mit der Stelle überein, an der Pallas nach der korrigierten Ephemeride stehen musste. Diese einzelne Beobachtung, die sich seither noch nicht hatte bestätigen lassen, gab aber noch keine volle Gewissheit, ob es wirklich Pallas gewesen war.\par}

\par\indent An der hiesigen Sternwarte ließ sich in der Nacht vom 4. auf den 5. Januar erstmals seit beinahe zwei Monaten wieder versuchen, Pallas zu finden. Da Professor \textsc{Gauss} dazu nur bewegliche Instrumente verwenden konnte, erschwerten der Mondschein und die strenge Kälte (11 Grad) die Suche ungemein. Seine Anstrengungen waren dennoch von Erfolg gekrönt. Professor \textsc{Gauss} fand Pallas an der Stelle wieder, an der er sie nach seiner korrigierten Ephemeride erwartet hatte. Außerdem gelangen ihm ein paar Vergleiche mit Sternen der Histoire céleste. Daraus leitete er unter dem Vorbehalt einer späteren genaueren Reduktion die folgende Position des Planeten ab.

\begin{center}\fontsize{10}{12}\selectfont\setlength{\tabcolsep}{4pt}\renewcommand{\arraystretch}{1.15}\begin{tabular}{c|c|c}
1811 Mittl. Zeit & Rektaszension & südl. Deklination\\\hline
Jan. 4 $15^{\text{u}}\;4^{\text{m}}\;42^{\text{s}}$ & $151^\circ\;40'\;5''$ & $22^\circ\;22'\;36''$\\
\end{tabular}\end{center}

\par\indent Der Planet hatte übrigens erst die 8. bis 9. Größe.

\Needspace{17\baselineskip}{\rightskip=0pt plus 1em \par\indent Auswärtige Astronomen werden den Planeten mit der eben angegebenen Korrektur der Ephemeride leicht finden. Wer die Ephemeride weiterführen oder genauer berechnen möchte, kann dazu die folgenden Elemente verwenden, die die gegenwärtige Bewegung des Planeten darstellen.\par}

\begin{center}\fontsize{10}{12}\selectfont\setlength{\tabcolsep}{4pt}\renewcommand{\arraystretch}{1.15}\begin{tabular}{llr}
\shortstack[l]{Epoche der mittleren Länge\\im Meridian von Göttingen 1811} & .\;. & $126^\circ\;33'\;6''$\\
Länge des Perihels & .\;.\;.\;.\;.\;.\;.\;.\;.\;.\;.\;.\;.\;.\;.\;. & $120\;55\;8$\\
Länge des aufsteigenden Knotens & .\;.\;.\;.\;.\;.\;.\;.\;.\;.\;.\;. & $172\;32\;40$\\
Neigung der Bahn & .\;.\;.\;.\;.\;.\;.\;.\;.\;.\;.\;.\;.\;.\;.\;.\;. & $34\;35\;14$\\
Tägliche mittlere tropische Bewegung & .\;.\;.\;.\;.\;.\;.\;,\;.\;. & $768''980$\\
Exzentrizität & .\;.\;.\;.\;.\;.\;.\;.\;.\;.\;.\;.\;.\;.\;.\;.\;.\;.\;.\;. & $0.24160$\\
Logarithmus der großen Halbachse & .\;.\;.\;.\;.\;.\;.\;.\;.\;.\;.\;. & $0.44278$\\
\end{tabular}\end{center}

\par\indent Professor \textsc{Gauss} glaubt, den praktischen Astronomen, die Pallas beobachten wollen, mit der unverzüglichen Bekanntgabe dieser Ergebnisse einen Dienst zu erweisen. Eine vollständige Darstellung seiner Arbeit über die Störungen der Pallas behält er sich für eine andere Gelegenheit vor.
